\begin{proposition}[Deformations of the constructed torus fibrations]\hypertarget{tag-01GW}{}\label{cdp:st:E7}
\StatementRecord{cdp:st:E7}{proved}{01GW}{\SeriesRef{cdp:st:E6}, \SeriesRef{cdp:eq:E-log-period}, \SeriesRef{cdp:st:E1}; \cite[Proposition~6.1]{Viaclovsky2026Combs}; \cite[Proposition~3.1]{Viaclovsky2026Two}; \cite[Proposition~2.5]{LiuRollenske2026}; \cite[Proposition~4.1, case~(i)]{Viaclovsky2026Combs}}

Every member of the following torus-fibred families,
including the period family $X_c$, its members with $|p|>1$,
Viaclovsky's $X_{3k,4l}$, his two-parameter family and his combs,
and the Liu--Rollenske family, satisfies the rank hypothesis of \ref{cdp:st:E6}.
Consequently every sufficiently small deformation retains a
holomorphic torus fibration over $\mathbb P^1$. Every such member
with $b_2=0$, and every sufficiently small deformation of it, has
algebraic dimension one. In particular the fibration-retention
question in \cite[Section~10(1)]{Viaclovsky2026Combs} has an
affirmative answer.
\end{proposition}


\begin{proof}

\begin{enumerate}[label=\arabic*.,leftmargin=*]
\item At their hexagonal fibre, a marking on the punctured disc puts
the Hodge periods in the form
\begin{equation}\hypertarget{tag-019Z}{}\label{cdp:eq:E-log-period}
  (I_2,Z(q)),\qquad
 Z(q)=\frac{\log q}{2\pi i}B+H(q),
 \qquad \det B\ne0,
\end{equation}
where $H$ is holomorphic at zero. Therefore
\[
 qZ'(q)=\frac{B}{2\pi i}+qH'(q)
\]
is invertible for all sufficiently small nonzero $q$. The
Kodaira--Spencer tensor of a marked torus with periods $(I,Z)$
is represented by $Z'$ after the invertible identifications of
the Lie algebra and $H^1(\mathcal O)$ with the chosen period
coordinates. Thus it has rank two at those regular fibres.

\item The local period models occur in
\cite[Proposition~6.1]{Viaclovsky2026Combs},
\cite[Proposition~3.1]{Viaclovsky2026Two}, and
\cite[Proposition~2.5]{LiuRollenske2026}.
The AC cusp has logarithmic matrix
$B_0=\left(\begin{smallmatrix}0&1\\-1&0\end{smallmatrix}\right)$.
The affine twists and finite logarithmic transformations change
the finite collars and constant gluing data, while the punctured
cusp matrix still has this same invertible logarithmic coefficient.
Engel's two deck transformations at the $I_1$-fibre have independent
valuation vectors $(1,0)$ and $(0,1)$, again giving an invertible
$B$. In the two-parameter family the designated fibre over $p_3$
is the full-lattice rank-two filling; its displayed two valuation
vectors are independent. Liu--Rollenske's compactification over
$I_1$ uses the same two independent valuation directions. Under the covering
$(z_1,z_2)\mapsto(e^{2\pi iz_1},e^{2\pi iz_2})$, lifts of the
angular loops give the two constant period columns, and logarithms
of the two deck multipliers, divided by $2\pi i$, give the columns of $Z(q)$.

\item For the combs, \cite[Proposition~4.1, case~(i)]{Viaclovsky2026Combs}
uses the flat cohomology basis $\alpha_i,\beta_i$ and the
single-valued holomorphic frame
$e_i=\alpha_i+\ell(q)\beta_i$, $f_i=\beta_i$, with
$\ell(q)=\log q/(2\pi i)$. The extended Hodge subbundle has
the rows $(I,H(q))$ in the frame $(e_i,f_i)$. The fixed value of the
smoothed period section at $q=0$ keeps it transverse to the
$f_i$-plane. Its rows therefore remain $(I,H_s(q))$ with
$H_s$ holomorphic. In the flat basis $(\alpha_i,\beta_i)$ the rows are
$(I,\ell(q)I+H_s(q))$, which is \eqref{cdp:eq:E-log-period} with
$B=I$.

\item All the sources construct the indicated members as compact
connected-fibre holomorphic maps with smooth general two-torus
fibre. \ref{cdp:st:E6} applies, and \ref{cdp:st:E1} gives the algebraic-dimension assertion.
On the differentiable six-sphere these neighbourhoods of the
listed structures constitute an open set on which $a=1$.

\end{enumerate}
\end{proof}